%% kapa: exam=1 round=26 year=2015 period=2 form=A subject=appraisal_law subject_ko=부동산관계법규 seq=30 number=30
%% source: Q-Net 감정평가사 기출문제 (cst00309) artlSeq=5138001 file=extracted/1008/5138001/제26회_감정평가사 1차_기출/제26회_감정평가사 1차_기출/2교시_감평_A형.pdf
%% conversion: kordoc 4.24.51 -> tools/gen1.py (auto); review=proofread-sonnet; flags=none
\begin{question}{30}
건축법령상 건축물의 건축허가를 받으면 허가를 받거나 신고를 한 것으로 본다. 이러한 허가 등의 의제에 해당하지 않는 것은?
\begin{choices}
  \item 「농지법」에 따른 농지전용허가
  \item 「하천법」에 따른 하천점용허가
  \item 「폐기물처리법」에 따른 폐기물처리업허가
  \item 「대기환경보전법」에 따른 대기오염물질 배출시설 설치의 신고
  \item 「국토의 계획 및 이용에 관한 법률」에 따른 개발행위허가
\end{choices}
\end{question}
\answer{3}
